\section{Related work}

\paragraph{Selective non-response as a confound.}
A cluster of recent work converges on the concern that a model's willingness
to answer is itself a confound in fairness evaluation: desirability-matched
graded forced choice with IRT-estimated latent traits
\citep{okada2026desirability}, clinical-style validity scales for
metacognitive self-report \citep{cacioli2026validity}, an IRT-fitted
``refusal strictness'' factor across 192 safety benchmarks
\citep{fonsecarivera2026irt}, implicit association testing
\citep{bai2025implicit}, list-experiment elicitation of concealed preferences
\citep{chupilkin2026hidden}, guardrail-agnostic probing via image-conveyed
demographics \citep{hirota2026guardrail}, activation steering
\citep{himelstein2026silenced}, and large correspondence-audit designs
\citep{gao2026hire}. Closest to our concern, \citet{gan2026biaxisbias} find
17.1\% of model--statement pairs change their forced-choice selection under
five equivalent wordings alone. UnQover \citep{li2020unqover} established the
forced-choice-only baseline this literature departs from, and
\citet{manduru2026vacuous} name the failure mode ``vacuous neutrality.'' This
body of work establishes the phenomenon. What has been missing is how much of
the published record it accounts for, and whether the natural remedy is
sound. We supply both.

\paragraph{Partial identification.} Given a non-response rate $p$ and an
answered-only alignment $q$, the sharp interval for population alignment is
$[(1-p)q,\ (1-p)q+p]$, assuming disclosed answers reveal the target outcome
and leaving undisclosed outcomes unrestricted \citep{manski1989anatomy}.
Interpreting this target as forced-choice behaviour additionally requires
that forcing preserve answers on disclosed cases.
\citet{horowitz1995identification} extend this to contaminated rather than
merely missing data, the closer analogy for selective disclosure. We use
these bounds as the identification backbone of \S\ref{sec:decomposition}.

\paragraph{Interpretability as an instrument.} Refusal in open-weight chat
models is mediated by a single residual-stream direction
\citep{arditi2024refusal}, which we use to build a refusal-ablated forced
condition, testing rather than assuming its validity: \S\ref{sec:models}
shows the certificate detecting where it fails and that conclusions hold when
it is replaced.

\paragraph{Evaluation data.} We use BBQ \citep{parrish2022bbq} for
\S\ref{sec:decomposition} and WinoGender
\citep{rudinger2018gender,rudinger2018winogender} for \S\ref{sec:models},
noting the template defects catalogued by WinoPron
\citep{gautam2024winopron}.
